\documentclass[tikz,border=4pt]{standalone}
\usepackage{amsmath,amssymb}
\usetikzlibrary{arrows.meta,calc}

% ============================================================================
% Adjustable model geometry (angles in degrees, lengths in cm)
% ============================================================================
\def\WheelCenterX{0.0}       % wheel-center coordinate O_w
\def\WheelCenterY{0.0}
\def\WheelRadius{1.05}       % wheel radius r
\def\WheelHubRadius{0.16}

\def\LegLength{3.65}         % exact distance |O_w H| = L
\def\ThetaDeg{15.0}          % theta: clockwise from global +y
\def\LegLineWidth{2.40pt}    % approximately 25 percent thinner than before

\def\HipJointRadius{0.18}
\def\BodyWidth{3.35}         % w_b, parallel to body longitudinal axis
\def\BodyHeight{1.20}        % h_b, normal to body longitudinal axis
\def\PhiDeg{12.0}            % phi: counterclockwise from global +x

% ============================================================================
% Adjustable annotation geometry
% ============================================================================
\def\XAxisLength{4.65}
\def\YAxisLength{5.35}
\def\GroundHalfLength{1.70}

\def\LocalVerticalStart{0.27}
\def\LocalVerticalLength{1.34}
\def\ThetaArcRadius{1.15}
\def\ThetaArcEndGapDeg{2.0}
\def\ThetaLabelGap{0.32}
\def\ThetaLabelAngularOffset{0.0}

\def\PhiGuideLength{1.42}
\def\PhiArcRadius{0.55}
\def\PhiLabelGap{0.61}
\def\PhiLabelAngularOffset{0.0}

\def\BodyDimensionGap{0.29}
\def\BodyExtensionOvershoot{0.10}

\def\LegDimensionOffset{0.52}
\def\LegDimensionExtensionOvershoot{0.04}

\def\WheelTorqueRadius{1.02}
\def\HipTorqueRadius{0.45}
\def\TorqueLabelGap{0.22}
\def\WheelTorqueStartDeg{238}
\def\WheelTorqueSweepDeg{-108}
\def\WheelTorqueLabelDeg{182}
\def\HipTorqueStartDeg{205}
\def\HipTorqueSweepDeg{-128}
\def\HipTorqueLabelDeg{150}

\def\JointLabelSep{7pt}
\def\HipLabelRadius{0.60}
\def\HipLabelAngle{225}
\def\DimensionLabelSep{2pt}

% ============================================================================
% Unified line-width hierarchy
% ============================================================================
\def\OutlineLineWidth{1.05pt}
\def\AxisLineWidth{0.90pt}
\def\DimensionLineWidth{0.48pt}
\def\ExtensionLineWidth{0.38pt}
\def\GuideLineWidth{0.55pt}
\def\TorqueLineWidth{0.95pt}

% ============================================================================
% Derived geometry: changing O_w, L, theta, r, phi, or body dimensions updates
% every dependent object automatically.
% ============================================================================
\pgfmathsetmacro{\HipJointX}{\WheelCenterX + \LegLength*sin(\ThetaDeg)}
\pgfmathsetmacro{\HipJointY}{\WheelCenterY + \LegLength*cos(\ThetaDeg)}

\pgfmathsetmacro{\LegNormalX}{cos(\ThetaDeg)}
\pgfmathsetmacro{\LegNormalY}{-sin(\ThetaDeg)}

\pgfmathsetmacro{\ThetaArcDelta}{-\ThetaDeg + \ThetaArcEndGapDeg}
\pgfmathsetmacro{\ThetaLabelRadius}{\ThetaArcRadius + \ThetaLabelGap}
\pgfmathsetmacro{\ThetaLabelAngle}{90 - \ThetaDeg/2 + \ThetaLabelAngularOffset}

\pgfmathsetmacro{\PhiLabelRadius}{\PhiArcRadius + \PhiLabelGap}
\pgfmathsetmacro{\PhiLabelAngle}{\PhiDeg/2 + \PhiLabelAngularOffset}

\pgfmathsetmacro{\WheelTorqueLabelRadius}{\WheelTorqueRadius + \TorqueLabelGap}
\pgfmathsetmacro{\HipTorqueLabelRadius}{\HipTorqueRadius + \TorqueLabelGap}

\begin{document}
\begin{tikzpicture}[
  x=1cm,
  y=1cm,
  line cap=round,
  line join=round,
  mechanism outline/.style={black,line width=\OutlineLineWidth},
  axis/.style={line width=\AxisLineWidth,
    -{Stealth[length=2.7mm,width=1.9mm]}},
  local guide/.style={blue!65!black,line width=\GuideLineWidth,
    densely dashed},
  body guide/.style={red!75!black,line width=\GuideLineWidth,
    densely dashed},
  angle arrow/.style={line width=0.82pt,
    -{Stealth[length=1.65mm,width=1.15mm]}},
  torque arrow/.style={orange!85!black,line width=\TorqueLineWidth,
    -{Stealth[length=2.4mm,width=1.75mm]}},
  dimension/.style={<->,>={Stealth[length=1.65mm,width=1.15mm]},
    line width=\DimensionLineWidth},
  extension/.style={line width=\ExtensionLineWidth},
  math label/.style={font=\large},
  angle label/.style={font=\small},
  dimension label/.style={font=\normalsize},
  torque label/.style={font=\normalsize},
  axis label/.style={font=\normalsize}
]

% ----------------------------------------------------------------------------
% Named model coordinates
% ----------------------------------------------------------------------------
\coordinate (Ow) at (\WheelCenterX,\WheelCenterY);
\coordinate (H) at (\HipJointX,\HipJointY);
\coordinate (GroundContact) at ($(Ow)+(0,-\WheelRadius)$);

% Leg-length dimension coordinates: offset from the leg along its right normal.
\coordinate (LegDimOw) at
  ($(Ow)+({\LegDimensionOffset*\LegNormalX},
           {\LegDimensionOffset*\LegNormalY})$);
\coordinate (LegDimH) at
  ($(H)+({\LegDimensionOffset*\LegNormalX},
          {\LegDimensionOffset*\LegNormalY})$);
\coordinate (LegExtOw) at
  ($(Ow)+({(\LegDimensionOffset+\LegDimensionExtensionOvershoot)*\LegNormalX},
           {(\LegDimensionOffset+\LegDimensionExtensionOvershoot)*\LegNormalY})$);
\coordinate (LegExtH) at
  ($(H)+({(\LegDimensionOffset+\LegDimensionExtensionOvershoot)*\LegNormalX},
          {(\LegDimensionOffset+\LegDimensionExtensionOvershoot)*\LegNormalY})$);

% Body-axis guide endpoint.
\coordinate (BodyAxisGuideEnd) at
  ($(H)+({\PhiGuideLength*cos(\PhiDeg)},
          {\PhiGuideLength*sin(\PhiDeg)})$);

% ----------------------------------------------------------------------------
% Ground: its height is derived from O_w and r, so it is tangent to the wheel.
% ----------------------------------------------------------------------------
\draw[gray!55,line width=\ExtensionLineWidth]
  ($(GroundContact)+(-\GroundHalfLength,0)$) --
  ($(GroundContact)+( \GroundHalfLength,0)$);

% ----------------------------------------------------------------------------
% Bottom layer: body first, then wheel. The body is centered at H and rotated
% by phi, so its long edge is exactly parallel to the body-axis guide.
% ----------------------------------------------------------------------------
\begin{scope}[shift={(H)},rotate=\PhiDeg]
  \draw[mechanism outline,fill=gray!8]
    ({-\BodyWidth/2},{-\BodyHeight/2}) rectangle
    ({ \BodyWidth/2},{ \BodyHeight/2});
\end{scope}

\draw[mechanism outline,fill=gray!4] (Ow) circle[radius=\WheelRadius];

% ----------------------------------------------------------------------------
% Middle layer: the leg centerline strictly joins O_w and H. Because it is
% drawn after the body, it remains visible all the way to the hip center.
% ----------------------------------------------------------------------------
\draw[black,line width=\LegLineWidth] (Ow) -- (H);

% ----------------------------------------------------------------------------
% Coordinate axes: drawn above the filled mechanism so they remain continuous.
% ----------------------------------------------------------------------------
\draw[axis,red!85!black]
  (Ow) -- ++(\XAxisLength,0)
  node[axis label,below right=1pt] {$x$};
\draw[axis,green!60!black]
  (Ow) -- ++(0,\YAxisLength)
  node[axis label,above left=1pt] {$y$};

% ----------------------------------------------------------------------------
% Local references and angle definitions
% ----------------------------------------------------------------------------
% Short local vertical reference at O_w. It starts outside the joint marker and
% is parallel to global +y.
\draw[local guide]
  ($(Ow)+(0,\LocalVerticalStart)$) --
  ($(Ow)+(0,\LocalVerticalLength)$);

% theta: clockwise from global +y toward the leg. The arc stops slightly before
% the leg centerline to prevent the arrowhead from overlapping the thick leg.
\draw[angle arrow,blue!70!black]
  ($(Ow)+(0,\ThetaArcRadius)$)
  arc[start angle=90,delta angle=\ThetaArcDelta,radius=\ThetaArcRadius];
\node[angle label,blue!70!black] at
  ($(Ow)+({\ThetaLabelRadius*cos(\ThetaLabelAngle)},
           {\ThetaLabelRadius*sin(\ThetaLabelAngle)})$)
  {$\theta$};

% phi: counterclockwise from the horizontal reference to the body long axis.
\draw[body guide] (H) -- ++(\PhiGuideLength,0);
\draw[body guide] (H) -- (BodyAxisGuideEnd);
\draw[angle arrow,red!80!black]
  ($(H)+(\PhiArcRadius,0)$)
  arc[start angle=0,delta angle=\PhiDeg,radius=\PhiArcRadius];
\node[angle label,red!80!black] at
  ($(H)+({\PhiLabelRadius*cos(\PhiLabelAngle)},
          {\PhiLabelRadius*sin(\PhiLabelAngle)})$)
  {$\phi$};

% ----------------------------------------------------------------------------
% Parametric dimensions and extension lines
% ----------------------------------------------------------------------------
% Leg length L = |O_w H|, offset parallel to the leg.
\draw[extension] (Ow) -- (LegExtOw);
\draw[extension] (H) -- (LegExtH);
\draw[dimension] (LegDimOw) -- (LegDimH)
  node[midway,right=\DimensionLabelSep,dimension label] {$L$};

% Body dimensions rotate with the body. Extension lines originate exactly at
% body corners and reach the corresponding parallel/perpendicular dimensions.
\begin{scope}[shift={(H)},rotate=\PhiDeg]
  % Width w_b, parallel to the body long edge.
  \draw[extension]
    ({-\BodyWidth/2},{\BodyHeight/2}) --
    ({-\BodyWidth/2},{\BodyHeight/2+\BodyDimensionGap+\BodyExtensionOvershoot});
  \draw[extension]
    ({ \BodyWidth/2},{\BodyHeight/2}) --
    ({ \BodyWidth/2},{\BodyHeight/2+\BodyDimensionGap+\BodyExtensionOvershoot});
  \draw[dimension]
    ({-\BodyWidth/2},{\BodyHeight/2+\BodyDimensionGap}) --
    ({ \BodyWidth/2},{\BodyHeight/2+\BodyDimensionGap})
    node[midway,above=\DimensionLabelSep,dimension label] {$w_b$};

  % Height h_b, perpendicular to the body long edge.
  \draw[extension]
    ({\BodyWidth/2},{-\BodyHeight/2}) --
    ({\BodyWidth/2+\BodyDimensionGap+\BodyExtensionOvershoot},{-\BodyHeight/2});
  \draw[extension]
    ({\BodyWidth/2},{ \BodyHeight/2}) --
    ({\BodyWidth/2+\BodyDimensionGap+\BodyExtensionOvershoot},{ \BodyHeight/2});
  \draw[dimension]
    ({\BodyWidth/2+\BodyDimensionGap},{-\BodyHeight/2}) --
    ({\BodyWidth/2+\BodyDimensionGap},{ \BodyHeight/2})
    node[midway,right=\DimensionLabelSep,dimension label] {$h_b$};
\end{scope}

% ----------------------------------------------------------------------------
% Actuator torques: common color, line width, arrow tip, and label gap.
% Both positive torque directions are clockwise on their target bodies.
% ----------------------------------------------------------------------------
\draw[torque arrow]
  ($(Ow)+({\WheelTorqueRadius*cos(\WheelTorqueStartDeg)},
           {\WheelTorqueRadius*sin(\WheelTorqueStartDeg)})$)
  arc[start angle=\WheelTorqueStartDeg,
      delta angle=\WheelTorqueSweepDeg,
      radius=\WheelTorqueRadius];
\node[torque label,orange!85!black] at
  ($(Ow)+({\WheelTorqueLabelRadius*cos(\WheelTorqueLabelDeg)},
           {\WheelTorqueLabelRadius*sin(\WheelTorqueLabelDeg)})$)
  {$T_w$};

\draw[torque arrow]
  ($(H)+({\HipTorqueRadius*cos(\HipTorqueStartDeg)},
          {\HipTorqueRadius*sin(\HipTorqueStartDeg)})$)
  arc[start angle=\HipTorqueStartDeg,
      delta angle=\HipTorqueSweepDeg,
      radius=\HipTorqueRadius];
\node[torque label,orange!85!black] at
  ($(H)+({\HipTorqueLabelRadius*cos(\HipTorqueLabelDeg)},
          {\HipTorqueLabelRadius*sin(\HipTorqueLabelDeg)})$)
  {$T_h$};

% ----------------------------------------------------------------------------
% Top layer: joint circles and their labels. These cover line endpoints cleanly
% while preserving the exact underlying O_w--H connection.
% ----------------------------------------------------------------------------
\draw[white,line width=2.8pt] (Ow) circle[radius=\WheelHubRadius];
\draw[mechanism outline,fill=white] (Ow) circle[radius=\WheelHubRadius];
\fill[black] (Ow) circle[radius=0.060];

\draw[white,line width=2.8pt] (H) circle[radius=\HipJointRadius];
\draw[mechanism outline,fill=white] (H) circle[radius=\HipJointRadius];
\fill[black] (H) circle[radius=0.060];

\node[math label,below right=\JointLabelSep] at (Ow) {$O_w$};
\node[math label] at
  ($(H)+({\HipLabelRadius*cos(\HipLabelAngle)},
          {\HipLabelRadius*sin(\HipLabelAngle)})$)
  {$H$};

\end{tikzpicture}
\end{document}
